% Part: first-order-logic
% Chapter: lindstrom
% Section: ls-property

\documentclass[../../../include/open-logic-section]{subfiles}

\begin{document}

\olfileid{mod}{lin}{lsp}

\olsection{సంహతత్వ, లొవెన్‌హైమ్--స్కోలెమ్ ధర్మాలు}

ఇప్పుడు నైరూప్య తర్కాల సందర్భానికి సంహతత్వం, లొవెన్‌హైమ్--స్కోలెమ్
ధర్మాల సహజ విస్తరణలను ఇస్తాం.

\begin{defn}
నైరూప్య తర్కం $\tuple{L,\models_L}$కు $L(\Lang{L})$-\tetoken{వాక్యాల}{!!{sentence}s} ప్రతి సమితి
$\Gamma$కు చెందిన ప్రతి పరిమిత $\Gamma_0\subseteq\Gamma$
సంతృప్తిపరచదగినదైతే ఆ సమితి కూడా సంతృప్తిపరచదగినదే అయితే,
\emph{సంహతత్వ ధర్మం} ఉంది అంటాం.
\end{defn}

\begin{defn}
సంతృప్తిపరచదగిన ప్రతి $\Gamma$కు
\tetoken{ఒక లెక్కించదగిన}{!!a{enumerable}} నమూనా ఉంటే,
$\tuple{L,\models_L}$కు \emph{అధోముఖ లొవెన్‌హైమ్--స్కోలెమ్ ధర్మం}
ఉంది అంటాం.
\end{defn}

\olref[bas][pis]{defn:partialisom}లోని పాక్షిక సమరూపత భావన పూర్తిగా
“బీజగణిత సంబంధమైనది” (అంటే, \tetoken{వాక్యాల}{!!{sentence}s}
ప్రస్తావన లేకుండా, \tetoken{భాష}{!!{language}}~$\Lang{L}$ అందించే
స్థిరసంకేతాల ఆధారంగా మాత్రమే ఇవ్వబడింది; ఆ నిర్మాణాలకే
\tetoken{నిర్మాణాలు}{!!{structure}s} చెందినవి); అందువల్ల అది
నైరూప్య తర్కాల సందర్భానికీ వర్తిస్తుంది. మొదటిస్థాయి తర్కంలో రెండు
\tetoken{నిర్మాణాలు}{!!{structure}s} పాక్షికంగా సమరూపమైతే అవి
మొదటిస్థాయి-వాక్య తుల్యాలని \olref[bas][pis]{thm:p-isom2} నుంచి
తెలుసు. \tetoken{సూత్రాలపై}{!!{formula}s} ఆగమనం యథేచ్ఛ
$!E\in L(\Lang{L})$కు అందుబాటులో ఉండనవసరం కాబట్టి, ఆ నిరూపణ
నైరూప్య తర్కాలకు వర్తించదు. అయినప్పటికీ లొవెన్‌హైమ్--స్కోలెమ్
ధర్మం ఉంటే సిద్ధాంతం సత్యమే.

\begin{thm}
\ollabel{thm:abstract-p-isom}
$\tuple{L,\models_L}$ అనేది లొవెన్‌హైమ్--స్కోలెమ్ ధర్మం గల సాధారణ
తర్కం అనుకుందాం. అప్పుడు పాక్షికంగా సమరూపమైన ఏ రెండు
\tetoken{నిర్మాణాలు}{!!{structure}s} అయినా $\tuple{L,\models_L}$లో
  వాక్య తుల్యాలు.
\end{thm}

\begin{proof}
$\Struct{M}\simeq_p\Struct{N}$ అయినప్పటికీ ఏదో $!E$కు
$\Struct{M}\models_L!E$, $\Struct{N}\not\models_L!E$ అనుకుందాం.
సమరూపతా ధర్మం వల్ల $\Domain{M}$, $\Domain{N}$లు పరస్పరం వియుక్తమని,
విస్తరణ ధర్మం వల్ల $!E\in L(\Lang{L})$ అయ్యే
\tetoken{భాష}{!!{language}}~$\Lang{L}$ పరిమితమని అనుకోవచ్చు.
$\Struct{M}$, $\Struct{N}$ల మధ్య పాక్షిక సమరూపతల సమితి
$\mathcal{I}$ అనుకుందాం. సాధారణత్వాన్ని కోల్పోకుండా, $p\in\mathcal{I}$,
$q\subseteq p$ అయితే $q\in\mathcal{I}$ కూడా అని అనుకుందాం.

$\Domain{M}^{<\omega}$ అనేది $\Domain{M}$లోని
\tetoken{మూలకాల}{!!{element}s} పరిమిత అనుక్రమాల సమితి. $S$ అనేది
$\Domain{M}^{<\omega}$పై అనుసంధానాన్ని సూచించే త్రిస్థాన సంబంధం
అనుకుందాం; అంటే $\mathbf{a},\mathbf{b},\mathbf{c}\in
\Domain{M}^{<\omega}$ అయితే, $\mathbf{c}$ అనేది $\mathbf{a}$,
$\mathbf{b}$ల అనుసంధానం అయినప్పుడు, అప్పుడు మాత్రమే
$S(\mathbf{a},\mathbf{b},\mathbf{c})$. అలాగే $b\in M$,
$\mathbf{a},\mathbf{c}\in\Domain{M}^{<\omega}$లకు,
  $\mathbf{a}=a_1, \dots a_n$, $\mathbf{c}=a_1, \dots a_n,b$ అయినప్పుడు,
అప్పుడు మాత్రమే $T(\mathbf{a},b,\mathbf{c})$ అయ్యే త్రిస్థాన సంబంధం
$T$ అనుకుందాం. కొత్త 3-స్థానాల \tetoken{విధేయ సంకేతాలు}{!!{predicate}s}
$P$, $Q$లను ఎంచుకొని, వ్యక్తి సముదాయం $\Domain{M}\cup
\Domain{M}^{<\omega}$గా ఉండి, $\Struct{M}$ను ఉపనిర్మాణంగా కలిగి, $P$,
$Q$లకు వరుసగా అనుసంధాన సంబంధాలు $S$, $T$లను అర్థనిర్దేశాలుగా ఇచ్చే
\tetoken{నిర్మాణం}{!!{structure}}~$\Struct{M}^*$ను నిర్మించండి
(కాబట్టి $\Struct{M}^*$ యొక్క \tetoken{భాష}{!!{language}}
$\Lang{L}\cup\{P,Q\}$).

$\Domain{N}^{<\omega}$, $S'$, $T'$, అలాగే అదే $P$, $Q$లకు వరుసగా
వాటిని అర్థనిర్దేశాలుగా ఇచ్చే $\Struct{N}^*$ను సమానంగా
నిర్వచించండి.
\sourcecorrection{OLTEMODLIN-017}{మొదటి వైపున అనుసంధానాన్ని
సంకేతీకరించడానికి కొత్త విధేయ సంకేతాలు $P,Q$ను ఎంచుకున్న తరువాత,
రెండవ వైపున ఆంగ్ల మూలం $P',Q'$ అనే వేరే సంకేతాలను జాబితా చేస్తుంది.
రెండు వైపులనూ ఒకే పరిసర భాషలో పోల్చి సాపేక్షీకరించాలంటే ఒకే
$P,Q$లను వరుసగా $S,T$కు, $S',T'$కు అర్థనిర్దేశాలుగా ఇవ్వాలి;
అదే చేశాం.}
అవసరమైతే సమరూపమైన ట్యాగ్‌చేసిన ప్రతులను తీసుకొని, అసలు మూలకాల
రెండు సముదాయాలు, వాటి పరిమిత-అనుక్రమ సముదాయాలు అన్నీ ఉద్దేశించిన
విధంగా పరస్పరం వియుక్తంగా ఉండేలా ఏర్పరచండి.
\sourcecorrection{OLTEMODLIN-021}{మూలకాల సముదాయాన్ని దాని పరిమిత
అనుక్రమాల సముదాయంతో, తరువాత రెండు వైపుల సముదాయాలను పరస్పరం,
ఆంగ్ల మూలం నేరుగా సమ్మేళనం చేస్తుంది. ఒక మూలకం తానే అనుక్రమం
కావచ్చును; ఖాళీ అనుక్రమం రెండు వైపులా ఒకటే కావచ్చును. తరువాతి
ఏకస్థాన విధేయాలు ఈ వేరు రకాలను గుర్తిస్తాయి కాబట్టి, సమరూపతను
మార్చకుండా ట్యాగ్‌చేసిన వియుక్త ప్రతులను ముందే తీసుకున్నామని
స్పష్టం చేశాం.}
\olref[bas][pis]{defn:partialisom}లోని నిర్మాణ పద్ధతిని అనుసరించి,
పరికల్పన ప్రకారం $\Struct{M}\simeq_p\Struct{N}$ కాబట్టి,
$\Domain{M}^{<\omega}$, $\Domain{N}^{<\omega}$ల మధ్య $I$ అనే సంబంధం
ఉంది; $I(\mathbf{a},\mathbf{b})$ అవుతుంది అప్పుడే సమాన పొడవు గల
$\mathbf{a}$, $\mathbf{b}$ల అనురూప స్థానాలను పంపే నిర్దేశం ఎంచుకున్న
పాక్షిక సమరూపతల కుటుంబానికి చెందినప్పుడు.
\sourcecorrection{OLTEMODLIN-005}{ఆంగ్ల మూలం పరిమిత అనుక్రమాలు
$\mathbf{a},\mathbf{b}$ “సమరూపమై ముందుకు-వెనుకకు షరతును
నెరవేరుస్తాయి” అంటుంది; కానీ సమరూపత, ముందుకు-వెనుకకు ధర్మం
నిర్దేశాలకు, వాటి కుటుంబానికి వర్తిస్తాయి. ముందే ఎంచుకున్న
$\mathcal{I}$లో అనురూప-స్థాన నిర్దేశం సభ్యమనే ఖచ్చితమైన
సంకేతీకరణను వాడాం.}
ఇప్పుడు \tetoken{నిర్మాణం}{!!{structure}}~$\Struct{K}$ యొక్క
\tetoken{వ్యక్తి క్షేత్రం}{!!{domain}}, $\Struct{M}^*$,
$\Struct{N}^*$ల \tetoken{వ్యక్తి క్షేత్రాల}{!!{domain}s}
సమ్మేళనం అనుకుందాం; $\Struct{K}$లో $\Struct{M}^*$,
$\Struct{N}^*$లు ఉప\-\tetoken{నిర్మాణాలు}{!!{structure}s}. దాని
\tetoken{భాషలో}{!!{language}} మరో ద్విస్థాన
\tetoken{విధేయ సంకేతం}{!!{predicate}}~$R$కు $I$ను అర్థనిర్దేశంగా
ఇచ్చి, \tetoken{విధేయ సంకేతాలను}{!!{predicate}s} కూడా చేర్చి,
వాటితో \tetoken{వ్యక్తి క్షేత్రాలు}{!!{domain}s} $\Domain{M^*}$,
$\Domain{N^*}$లను సూచించండి.
\sourcecorrection{OLTEMODLIN-006}{అసలు నిర్మాణం $\Struct{M}$,
దాని విస్తరణ $\Struct{M}^*$ ఇప్పటికే ఉన్నప్పటికీ, ఆంగ్ల మూలం పరిసర
నిర్మాణానికీ $\Struct{M}$ అనే పేరే మళ్లీ పెట్టి, దానిలో
$\Struct{M}^*$తో పాటు $\Struct{M}$ తానే ఉపనిర్మాణమనే గందరగోళాన్ని
కలిగిస్తుంది; $\Domain{N}*$ను కూడా తప్పుగా నక్షత్రం వెలుపల
రాస్తుంది. పరిసర నిర్మాణాన్ని $\Struct{K}$గా, తరువాత దాని
లెక్కించదగిన నమూనాను $\Struct{K}_0$గా వేరు చేసి,
$\Domain{M^*}$, $\Domain{N^*}$లను సరిగ్గా రాశాం.}

\begin{figure}[h]
  \centering
  \begin{tikzpicture}[node distance=2cm, auto, thick, >=stealth']
    \draw [rounded corners] (0,0) -- (8,0) -- (8,4) -- (0,4) --  cycle;
    \draw (2,2) circle (0.5cm);
    \draw (2,2) circle (1.25cm);
    \draw (6,2) circle (0.5cm);
    \draw (6,2) circle (1.25cm);
    \path node at (0.75,3.5) {\large $\Struct{K}$};
    \path node at (2,2) {\large $\Struct{M}$};
    \path node at (6,2) {\large $\Struct{N}$};
    \path node at (3.5,1) {\large $\Struct{M}^*$};
    \path node at (7.5,1) {\large $\Struct{N}^*$};
    \node (Idom) at (2.8,2) {};
    \node (Irng) at (5.2,2) {};
    \draw[<->, bend left] (Idom) to node {\large $I$} (Irng) ;
  \end{tikzpicture}
  \caption{అంతర్గత పాక్షిక సమరూపత గల
    \tetoken{నిర్మాణం}{!!{structure}}~$\Struct{K}$.}
\end{figure}

కీలకమైన పరిశీలన ఇది: ఈ పరిసర \tetoken{భాషలో}{!!{language}},
\tetoken{నిర్మాణం}{!!{structure}}~$\Struct{K}$లో సత్యమైన,
$\Struct{M}\models_L!E$, $\Struct{N}\not\models_L!E$ అని చెప్పే
  నైరూప్య \tetoken{వాక్యం}{!!{sentence}}~$!D_1$ ఉంది (దీనికి
సాపేక్షీకరణ ధర్మం అవసరం). అలాగే, $\Struct{M}\simeq_p\Struct{N}$
అనేది పాక్షిక సమరూపత~$I$ ద్వారా నెరవేరుతుందని చెప్పే, $\Struct{K}$లో
సత్యమైన మొదటిస్థాయి \tetoken{వాక్యం}{!!{sentence}}~$!D_2$ ఉంది.
\sourcecorrection{OLTEMODLIN-007}{యథేచ్ఛ నైరూప్య $L$-వాక్యం
$!E$ యొక్క సత్యత, అసత్యతలను సాపేక్షీకరించి తెలిపే $!D_1$ను ఆంగ్ల
మూలం మొదటిస్థాయి వాక్యం అంటుంది. సాపేక్షీకరణ, బూలియన్ సంవృతత
దాన్ని $L$-వాక్యంగా ఇస్తాయి; పాక్షిక సమరూపతను సంకేతీకరించే
$!D_2$ మాత్రమే మొదటిస్థాయి వాక్యం. ఈ రకాలను వేరుచేశాం.}
సాధారణ తర్కపు బూలియన్ ధర్మం వల్ల $!D_1$, $!D_2$లను ఒకే
  నైరూప్య వాక్యంగా సంయోగించవచ్చు. లొవెన్‌హైమ్--స్కోలెమ్ ధర్మం వల్ల అది
\tetoken{ఒక లెక్కించదగిన}{!!a{enumerable}} నమూనా
$\Struct{K}_0$లో సత్యం. ఆ నమూనాలో $\Struct{M}_0$,
$\Struct{N}_0$ అనే పాక్షికంగా సమరూపమైన ఉపనిర్మాణాలు ఉండి,
$\Struct{M}_0\models_L!E$, $\Struct{N}_0\not\models_L!E$.
కానీ \tetoken{లెక్కించదగిన}{!!{enumerable}}, పాక్షికంగా సమరూపమైన
\tetoken{నిర్మాణాలు}{!!{structure}s} నిజానికి
\olref[bas][pis]{thm:p-isom1} ప్రకారం సమరూపాలు. ఇది సాధారణ నైరూప్య
తర్కాల సమరూపతా ధర్మానికి విరుద్ధం.
\end{proof}

\end{document}
